\section{Optimal Power Flow Models}
\label{sec:opf}

\subsection{Implemented OPF Surfaces and Model Spaces}

The repository currently has two OPF-facing formulation surfaces.  Both use the
canonical component model, but they serve different callers:
\begin{itemize}
  \item \texttt{hacdcpf::opf::solve\_ac\_opf} is the public production AC/hybrid
  OPF entry point.  It canonicalizes the rich \texttt{HybridPowerSystem} through
  \texttt{powerflow::make\_solver\_data}, which calls
  \texttt{project\_to\_canonical\_models}.  It then chooses either the
  full-space parity NLP, the older native reduced-space primal-dual model, or
  the pure-AC economic-dispatch fallback.  The parity NLP is the implemented
  path that optimizes AC, DC, VSC, and DC/DC dispatch together.
  \item \texttt{hacdcpf::power\_models::solve\_acopf},
  \texttt{solve\_acdcopf}, and \texttt{solve\_dc\_opf} are AML model builders.
  They expose solver-agnostic canonical OPF blocks for tests and direct AML
  callers.  These builders are simpler than the public parity OPF: they do not
  include load shedding, enhanced DER dispatch, storage, flexible loads,
  energy-router ports, or the public fallback chain.
\end{itemize}

This distinction is important for downstream ONR work.  Power flow and public
OPF operate on canonical branches after transformers, closed switches, and
other rich devices are projected.  Any ONR formulation that optimizes line
status should therefore define binary line-status decisions in the same
canonical branch space and keep a separate mapping back to the switching
devices or circuit breakers that realize that line status.

\subsection{Code Map}

\begin{itemize}
  \item \texttt{src/optimal\_power\_flow/ac\_opf.cpp}: public
  \texttt{solve\_ac\_opf} routing, economic-dispatch fallback, native
  reduced-space primal-dual solver, and Ipopt adapter for the parity NLP.
  \item \texttt{include/hacdcpf/optimal\_power\_flow/formulation.hpp} and
  \texttt{src/optimal\_power\_flow/parity\_formulation.cpp}: canonical
  full-space parity NLP assembly, objective, bounds, equality constraints,
  nonlinear inequalities, Jacobian, and Hessian.
  \item \texttt{src/optimal\_power\_flow/parity\_ipm.cpp}: native
  Mehrotra/Gondzio primal-dual IPM for the parity NLP.
  \item \texttt{src/optimal\_power\_flow/dc\_opf.cpp}: public DC OPF QP/LP
  model and backend selection.
  \item \texttt{src/power\_models/acopf\_builder.cpp},
  \texttt{src/power\_models/acdcopf\_builder.cpp}, and
  \texttt{src/power\_models/dc\_opf\_builder.cpp}: canonical AML OPF builders.
  \item \texttt{src/optimal\_power\_flow/reactive\_power\_opt.cpp}: discrete
  tap/shunt reactive-power optimization using repeated AC OPF evaluations.
\end{itemize}

\subsection{Canonicalization Contract}

The public AC/hybrid OPF path calls \texttt{powerflow::make\_solver\_data}.  That
routine first projects the rich input to canonical form:
\begin{align}
\widehat{\mathcal{S}}=\Pi_{\qual{canon}}(\mathcal{S}).
\end{align}
The projection expands rich devices into solver-visible canonical components and
stores recovery metadata such as bus-merge maps.  The resulting
\texttt{SolverData} contains flat arrays for AC buses and branches, DC buses and
branches, VSC converters, DCDC converters, energy routers, generators, loads,
DERs, storage, shunts, and other fixed injections.

The OPF equations below are written in this canonical space.  User-facing result
vectors that are per-AC-bus are expanded through the bus-merge map before being
returned where applicable.

\subsection{Public AC/Hybrid OPF Routing}

The public \texttt{solve\_ac\_opf} entry point first performs a graph topology
pre-check for usable AC islands.  It then selects the formulation from
\texttt{ACOPFOptions::ac\_solver\_backend} and the legacy
\texttt{enable\_primal\_dual}/\texttt{use\_parity\_ipm} flags.

\begin{longtable}{L{0.25\linewidth}L{0.28\linewidth}L{0.37\linewidth}}
\toprule
Path & Trigger & Implemented behavior \\
\midrule
\endhead
Parity full-space NLP &
\texttt{use\_parity\_ipm=true}; also auto-enabled for hybrid AC/DC cases unless
\texttt{EconomicDispatch} is explicitly requested &
Builds \texttt{opf::parity::Problem} from canonical \texttt{SolverData}; solves
with native parity IPM or Ipopt; includes DCDC active-power dispatch. \\
Native reduced-space primal-dual &
\texttt{enable\_primal\_dual=true}, \texttt{use\_parity\_ipm=false} &
Uses custom sparse KKT iterations over a smaller AC/DC/VSC variable set; DCDC
dispatch is not part of this older variable set. \\
Economic dispatch + AC PF &
\texttt{enable\_primal\_dual=false} on pure-AC cases &
Solves generator economic dispatch and projects by AC power flow.  Suppressed
for hybrid cases because it would drop DC, VSC, and DCDC equations. \\
\bottomrule
\end{longtable}

When \texttt{ACOPFSolverBackend::Auto} is used on a case with DC-side
components, including DCDC converters, the code enables the parity route.  If
native parity IPM fails and
\texttt{allow\_fallback=true}, the same assembled parity problem is retried with
Ipopt when Ipopt is compiled and not disabled by
\texttt{HACDCPF\_DISABLE\_IPOPT\_OPF}.

\subsection{Economic Dispatch + PF Fallback}

The pure-AC fast fallback solves:
\begin{align}
\min_{P_g}\quad
\sum_{g\in\mathcal{G}}\left(c_{2g}P_g^2+c_{1g}P_g+c_{0g}\right)
\end{align}
subject to generator bounds and a total active-power balance
\(\sum_g P_g=P_D\), clamped to the feasible generator range.  Quadratic units
use incremental-cost clipping,
\begin{align}
P_g=\operatorname{clip}\left(\frac{\lambda-c_{1g}}{2c_{2g}},
P_g^{\min},P_g^{\max}\right).
\end{align}
The dispatch is then passed to AC Newton power flow.  If the first projection
fails, one reshaped dispatch attempt is tried before failure is reported.

\subsection{Native Reduced-Space Primal-Dual OPF}

\subsubsection{Decision Variables}

The native path keeps only the non-slack AC angles, PQ-bus voltage magnitudes,
online generator powers, all DC-bus voltages, and VSC powers:
\begin{align}
x_{\qual{nat}}=
\begin{bmatrix}
\theta_{\mathcal{N}_{ac}\setminus\mathcal{N}_{sl}}&
V_{\mathcal{N}_{PQ}}&
P_g&Q_g&
V^{dc}&
P_c^{ac}&Q_c^{ac}&P_c^{dc}
\end{bmatrix}^{\top}.
\end{align}
Reference-bus angles and PV/slack voltage magnitudes are held at their canonical
setpoints.  DC voltage bounds are wide by default; \texttt{DC\_V} buses receive
a tight band around \texttt{vm\_pu}.

\subsubsection{Objective}

\begin{align}
\min_x\quad
f_{\qual{nat}}(x)=
\sum_{g\in\mathcal{G}}
\left(c_{2g}(S_bP_g)^2+c_{1g}S_bP_g+c_{0g}\right).
\end{align}

\subsubsection{Equality Constraints}

The native residual uses the existing power-flow residual kernel for AC network
balance and subtracts explicit converter variables.  With \(P_i^{\qual{PF}}\)
and \(Q_i^{\qual{PF}}\) denoting the PF specified-minus-calculated residuals
with the converter table disabled:
\begin{align}
g_{P,i} &=
P_i^{\qual{PF}}-\sum_{c\in\mathcal{C}^{ac}_i}P_c^{ac}=0,
&& i\in\mathcal{N}_{ac}\setminus\mathcal{N}_{sl},\\
g_{Q,i} &=
Q_i^{\qual{PF}}-\sum_{c\in\mathcal{C}^{ac}_i}Q_c^{ac}=0,
&& i\in\mathcal{N}_{PQ}.
\end{align}

The DC bus balance is:
\begin{align}
g_k^{dc}=
P_{d,k}^{dc}
 + V_k^{dc}\left(G^{dc}V^{dc}\right)_k
 - \sum_{c\in\mathcal{C}^{dc}_k}P_c^{dc}
=0 .
\end{align}

For each VSC converter:
\begin{align}
g_c^{\qual{vsc}}&=P_c^{ac}+P_c^{dc}+P_c^{loss}=0,\\
P_c^{loss}&=a_c+b_c I_c^{ac}+(1-\eta_c)(I_c^{ac})^2,\\
I_c^{ac}&=
\frac{\sqrt{(P_c^{ac})^2+(Q_c^{ac})^2+\epsilon}}{V_c^{ac}},
\qquad \epsilon=10^{-6}.
\end{align}

\subsubsection{Bounds and Nonlinear Engineering Limits}

The native path enforces variable bounds through a logarithmic barrier.  AC
branch thermal limits and VSC apparent-power limits are added to the merit model
as exterior quadratic penalties:
\begin{align}
(P_{f,\ell})^2+(Q_{f,\ell})^2-(S_{\ell}^{max})^2&\le 0,\\
(P_{t,\ell})^2+(Q_{t,\ell})^2-(S_{\ell}^{max})^2&\le 0,\\
(P_c^{ac})^2+(Q_c^{ac})^2-(S_c^{max})^2&\le 0.
\end{align}

At each iteration the native solver forms a regularized sparse KKT system:
\begin{align}
\begin{bmatrix}
H+\delta I & J_g^{\top}\\
J_g & -\delta I
\end{bmatrix}
\begin{bmatrix}
\Delta x\\
\Delta\lambda
\end{bmatrix}
=-
\begin{bmatrix}
\nabla_x L\\
g(x)
\end{bmatrix}.
\end{align}
Line search, trust/restore logic, bounded step clipping, and adaptive penalty
weights globalize the iterations.

\subsection{Parity Full-Space Canonical NLP}

\subsubsection{Decision Variables}

The parity formulation is the public hybrid-capable NLP.  It keeps all AC bus
angles and voltage magnitudes explicitly, followed by generation, converter,
shedding, enhanced DER, DCDC, flexible-load, and energy-router-port variables:
\begin{align}
x_{\qual{par}}=
\left[
\begin{array}{l}
\theta,\ V,\ P_g,\ Q_g,\ V^{dc},\ P_c^{ac},\ Q_c^{ac},\ P_c^{dc},\\
\Delta P_d,\ \Delta Q_d,\ P^{ren},\ Q^{ren},\ P^{stor},\ Q^{stor},\\
P^{stor,dc},\ P^{dcdc},\ P^{flex},\ P^{er},\ Q^{er}
\end{array}
\right]^{\top}.
\end{align}

Block lengths and offsets are defined by
\texttt{opf::parity::VarIndex}.  Empty component families have zero-length
blocks.  \(P^{dcdc}\) is an optimized DC/DC converter active-power dispatch
block, not a fixed post-processing quantity.

\subsubsection{Objective}

All power variables are stored in per-unit; cost terms are evaluated in
engineering units through \(S_b=\texttt{base\_mva}\):
\begin{align}
\min_x\quad
f_{\qual{par}}(x)=&
\sum_{g\in\mathcal{G}}
\left(c_{2g}(S_bP_g)^2+c_{1g}S_bP_g+c_{0g}\right) \nonumber\\
&+\mathrm{VOLL}_{eff}S_b
\sum_{i\in\mathcal{N}_{ac}}(\Delta P_{d,i}+\Delta Q_{d,i}) \nonumber\\
&+
\sum_{r\in\mathcal{R}^{curt}}
c_r^{curt}\max\left(0,P_r^{rated}-S_bP_r^{ren}\right).
\end{align}
\(\mathrm{VOLL}_{eff}\) is either supplied or auto-computed as at least
\(100\) and \(1.1\) times the maximum generator marginal cost.

\subsubsection{AC Balance Rows}

The AC network injection uses the polar Y-bus equations:
\begin{align}
P_i^{net}&=
V_i\sum_j V_j
\left(G_{ij}\cos\theta_{ij}+B_{ij}\sin\theta_{ij}\right),\\
Q_i^{net}&=
V_i\sum_j V_j
\left(G_{ij}\sin\theta_{ij}-B_{ij}\cos\theta_{ij}\right).
\end{align}

Voltage-dependent demand uses the same ZIP convention as the power-flow
assembly:
\begin{align}
P_i^d(V_i)&=\bar P_i^d(w_{P,i}^{P}+w_{I,i}^{P}V_i+w_{Z,i}^{P}V_i^2),\\
Q_i^d(V_i)&=\bar Q_i^d(w_{P,i}^{Q}+w_{I,i}^{Q}V_i+w_{Z,i}^{Q}V_i^2).
\end{align}

The implemented active and reactive balance rows are:
\begin{align}
g_{P,i}^{ac}=&
\alpha_P\Bigl(
P_i^{net}
-P_i^g-P_i^{fixed}-P_i^{ren}-P_i^{stor}
+P_i^d(V_i)-\Delta P_{d,i}
+P_i^{flex}
-P_i^{vsc,ac}
-P_i^{er}
\Bigr)=0,\\
g_{Q,i}^{ac}=&
\alpha_Q\Bigl(
Q_i^{net}
-Q_i^g-Q_i^{fixed}-Q_i^{ren}-Q_i^{stor}
+Q_i^d(V_i)-\Delta Q_{d,i}
-Q_i^{vsc,ac}
-Q_i^{er}
\Bigr)=0.
\end{align}
The scaling factors \(\alpha_P\) and \(\alpha_Q\) are reciprocal demand-scale
normalizers computed during problem assembly.  Fixed injections include
non-variable static generators, non-curtailable renewables/PV, VPPs,
grid-connected microgrid exchange, and mobile storage.

\subsubsection{DC Balance Rows}

For each canonical DC bus \(k\), the base conductance-flow balance is:
\begin{align}
g_k^{dc}=
P_{d,k}^{dc}
+\sum_{m\ne k}G_{km}^{dc}V_k^{dc}(V_k^{dc}-V_m^{dc})
-P_k^{vsc,dc}
+\cdots=0 .
\end{align}

The remaining terms are added directly to the same nodal rows:
\begin{itemize}
  \item DC loads increase demand.
  \item DC static generators are currently added with the sign implemented in
  the evaluator; DC PV arrays subtract their set power.
  \item DC storage contributes \(+P^{stor,dc}\).
  \item DC energy-router ports subtract \(P^{er}\) at their DC bus.
\end{itemize}

For a DCDC converter, the OPF variable \(P^{dcdc}\) is a true dispatch variable
optimized by the parity NLP.  It is input-side referenced: positive power is
drawn from \texttt{bus\_in}.  Its output-side power is a deterministic function,
\begin{align}
P^{out}=
\begin{cases}
\eta P^{in}-r_{eq}(P^{in})^2/(V_{in}^{dc})^2, & P^{in}\ge 0,\\
P^{in}/\eta-r_{eq}(P^{in})^2/(V_{in}^{dc})^2, & P^{in}<0.
\end{cases}
\end{align}
The contribution is folded into the two DC bus rows:
\begin{align}
g_{\texttt{bus\_in}}^{dc}&\mathrel{+}=P^{in},&
g_{\texttt{bus\_out}}^{dc}&\mathrel{-}=P^{out}.
\end{align}
There is no dedicated DCDC equality block in the current implementation
(\texttt{ConstraintIndex::n\_dcdc\_bal=0}).

DC buses with no conductive DC branch receive a soft voltage anchor
\(K(V_k^{dc}-V_{k,0}^{dc})\) in the DC balance row to avoid an unobservable
voltage column.

\subsubsection{VSC and Energy-Router Rows}

For each VSC converter:
\begin{align}
P_c^{ac}+P_c^{dc}+a_c+b_cI_c^{ac}+(1-\eta_c)(I_c^{ac})^2=0,
\end{align}
where
\begin{align}
I_c^{ac}=
\frac{\sqrt{(P_c^{ac})^2+(Q_c^{ac})^2+\epsilon}}{V_{a(c)}},
\qquad \epsilon=10^{-6}.
\end{align}

Each in-service energy router contributes one active-power conservation row:
\begin{align}
\sum_{p\in\mathcal{P}(r)}P_{r,p}^{er}=0.
\end{align}
AC ports also have reactive-power variables and enter AC \(Q\)-balance rows.
The current row is lossless; router loss refinements are not part of this OPF
formulation.

\subsubsection{Variable Bounds}

The parity bound builder enforces:
\begin{itemize}
  \item \(-\pi\le\theta_i\le\pi\) and AC/DC voltage limits from canonical buses;
  \item generator \(P/Q\) limits, with out-of-service generators fixed near their
  scheduled values;
  \item VSC \(P^{ac}/Q^{ac}\) limits and symmetric \(P^{dc}\) bounds inferred
  from converter capacity;
  \item load-shedding windows
  \(0\le\Delta P_{d,i}\le\max(P_{d,i},0.01/S_b)\) and similarly for \(Q\);
  \item curtailable renewable/PV availability and reactive limits;
  \item AC storage dispatch fixed to the scheduled active-power setpoint within
  a small band, with reactive limits;
  \item DC storage active-power limits;
  \item DCDC active-power limits from \((p_{\min},p_{\max})\) or rating data;
  \item flexible-load bounds from \(\texttt{flex\_down\_mw}\) and
  \(\texttt{flex\_up\_mw}\);
  \item energy-router port \(P/Q\) capability bounds.
\end{itemize}

\subsubsection{Nonlinear Inequalities}

Nonlinear inequalities use the convention \(h(x)\le 0\).  Row families are
counted by \texttt{ConstraintIndex}:
\begin{align}
(P_{f,\ell})^2+(Q_{f,\ell})^2-(S_\ell^{max})^2&\le0,\\
(P_{t,\ell})^2+(Q_{t,\ell})^2-(S_\ell^{max})^2&\le0,\\
(P_c^{ac})^2+(Q_c^{ac})^2-(S_c^{max})^2&\le0,\\
\left(G_{km}^{dc}V_k^{dc}(V_k^{dc}-V_m^{dc})\right)^2
-(P_{km}^{max})^2&\le0,\\
(P_c^{ac})^2+(Q_c^{ac})^2
-(I_{c,max}^{ac}V_{a(c)})^2&\le0.
\end{align}
Optional VSC modulation rows are linear in \(V^{ac}\) and \(V^{dc}\):
\begin{align}
V_a V_n^{ac}-m_{\max}K_mV_d^{dc}V_n^{dc}&\le0,\\
m_{\min}K_mV_d^{dc}V_n^{dc}-V_aV_n^{ac}&\le0.
\end{align}
Optional DCDC duty-ratio feasibility rows are also linear in the solved DC port
voltages after applying the topology-specific ideal conversion law.  The same
\texttt{enforce\_converter\_modulation\_limits} option gates VSC modulation and
DCDC duty-ratio rows.

The native parity IPM appends finite variable bounds as additional inequality
rows \(x^{min}-x\le0\) and \(x-x^{max}\le0\).  The Ipopt adapter instead passes
box bounds as solver variable bounds and only exposes the nonlinear rows above
as general inequalities.

\subsubsection{Parity IPM System}

The native parity solver applies a primal-dual interior-point method to:
\begin{align}
g(x)&=0,\\
h(x)+z&=0,\\
\nabla f(x)+J_g(x)^{\top}\lambda+J_h(x)^{\top}\mu&=0,\\
z\circ\mu&=\gamma\mathbf{1}.
\end{align}
It uses a Mehrotra predictor-corrector step, adaptive centering, optional
Gondzio extra correctors, fraction-to-boundary step lengths, sparse or dense
linear solves, iterative refinement, and best-iterate restoration before final
failure.

\subsection{Component Coverage}

\begin{longtable}{L{0.32\linewidth}L{0.26\linewidth}L{0.30\linewidth}}
\toprule
Component / feature & Native reduced-space path & Parity full-space path \\
\midrule
\endhead
Generator \(P_g,Q_g\) & decision variable & decision variable \\
AC bus voltage & non-slack angle and PQ magnitude variables & all angles and magnitudes \\
VSC converter \(P^{ac},Q^{ac},P^{dc}\) & decision variable & decision variable \\
DC bus voltage & decision variable & decision variable \\
AC branch MVA limits & exterior penalty in merit model & inequality rows when enabled \\
DC branch power limits & not explicit & inequality rows when enabled \\
VSC current/modulation limits & not explicit & optional inequality rows \\
Load shedding & not modeled & decision variable with VOLL penalty \\
Curtailable renewable / controllable PV & not modeled as variable & decision variable \\
AC storage dispatch & not modeled as variable & scheduled active dispatch with reactive variable \\
DC storage dispatch & not modeled as variable & active-power variable \\
DCDC dispatch & not modeled & optimized input-referenced variable folded into DC balances \\
Flexible load & not modeled & active-power variable \\
Energy-router ports & not modeled & port power variables plus router balance rows \\
\bottomrule
\end{longtable}

\subsection{AML Canonical OPF Builders}

The AML builders under \texttt{src/power\_models} are direct canonical model
builders and should not be confused with \texttt{opf::solve\_ac\_opf}.

\subsubsection{AML AC OPF}

\texttt{power\_models::to\_acopf\_data} keeps in-service canonical AC buses,
generators, and finite-impedance in-service AC branches.  Zero-impedance
switch/tie branches are skipped by this builder.  The decision variables are:
\begin{align}
x_{\qual{aml-ac}}=[V,\theta,P_g,Q_g]^{\top}.
\end{align}
The objective is:
\begin{align}
\min\sum_g
\left(c_{2g}(S_bP_g)^2+c_{1g}S_bP_g+c_{0g}\right).
\end{align}
For each AC bus, the builder accumulates:
\begin{align}
0=&
-P_{d,i}-G_{s,i}V_i^2+\sum_{g\in i}P_g
-\sum_{\ell\in\delta(i)}P_{\ell,i},\\
0=&
-Q_{d,i}+B_{s,i}V_i^2+\sum_{g\in i}Q_g
-\sum_{\ell\in\delta(i)}Q_{\ell,i}.
\end{align}
It pins the reference-bus angle by equality and adds both-end branch thermal
inequalities when \(\texttt{rate\_a\_pu}>0\).

\subsubsection{AML Hybrid AC/DC OPF}

\texttt{power\_models::solve\_acdcopf} extends the AML AC OPF with:
\begin{align}
x_{\qual{aml-acdc}}=
[V,\theta,P_g,Q_g,V^{dc},P_c^{ac},Q_c^{ac}]^{\top}.
\end{align}
Converter AC injections are added to AC bus balances with injection-positive
sign convention.  DC branch power is linear in the solved DC voltage variables:
\begin{align}
P_{km}^{dc}=\frac{V_k^{dc}-V_m^{dc}}{r_{km}}.
\end{align}
For each converter, the builder computes:
\begin{align}
P_c^{dc}=-P_c^{ac}-a_c-b_cI_c^{ac}-c_c(I_c^{ac})^2
\end{align}
inside the expression graph and adds it directly to the connected DC-bus
balance.  Thus, unlike the public parity formulation, this AML hybrid builder
has no explicit \(P_c^{dc}\) decision variable and no separate VSC equality row.
\texttt{VDC\_Q} converters add an equality \(V_k^{dc}=V_{set}^{dc}\).
This AML hybrid builder does not currently add DCDC dispatch variables; the
public parity formulation above is the implemented OPF path for optimized
DC/DC converter dispatch.

\subsubsection{AML DC OPF}

\texttt{power\_models::solve\_dc\_opf} is a linear AML DC OPF.  It uses:
\begin{align}
x_{\qual{aml-dc}}=[p_g,\theta]^{\top}
\end{align}
with generator bounds, a fixed slack angle, nodal balance, and branch flow box
limits:
\begin{align}
\sum_{g\in i}p_g-P_{d,i}
-\sum_{j:(i,j)}b_{ij}(\theta_i-\theta_j)
+\sum_{j:(j,i)}b_{ji}(\theta_j-\theta_i)&=0,\\
-F_{\ell}^{max}\le b_{\ell}(\theta_f-\theta_t)&\le F_{\ell}^{max}.
\end{align}
The objective is linear:
\begin{align}
\min\sum_g c_g p_g.
\end{align}

\subsection{Public DC OPF}
\label{sec:dc-opf}

The public \texttt{hacdcpf::opf::solve\_dc\_opf} path in
\texttt{src/optimal\_power\_flow/dc\_opf.cpp} is separate from the AML DC OPF
builder.  It constructs a DC-power-flow dispatch model with:
\begin{align}
x_{\qual{dc}}=
\begin{bmatrix}
\theta_{\mathcal{N}\setminus\{sl\}}&
P_g&
P_f&
\Delta P_d
\end{bmatrix}^{\top}.
\end{align}
Power balance and flow-definition rows are:
\begin{align}
\sum_{g\in\mathcal{G}_i}P_g-P_{d,i}+\Delta P_{d,i}
&=\sum_{j\in\mathcal{N}_i}P_{ij},\\
P_{ij}-b_{ij}(\theta_i-\theta_j)&=0.
\end{align}

The solver builds both a true QP objective for QP-capable backends and a
linearized LP objective for LP fallback:
\begin{align}
\min_{\qual{QP}}\quad&\frac{1}{2}x^{\top}Qx+c^{\top}x,\\
\min_{\qual{LP}}\quad&
\sum_g c_g^{lin}P_g+\sum_i\mathrm{VOLL}_{eff}S_b\Delta P_{d,i}.
\end{align}
The auto backend sequence tries \texttt{NativeQP}, Gurobi, HiGHS, and native
dual simplex.  \texttt{DCOPFResult::solver\_chain} records attempted backends,
and \texttt{objective\_model} records whether the reported objective used the
QP or LP model.

If \texttt{compute\_lmp=true}, a supporting simplex LP is run after the primal
solve to recover nodal prices.  Branch shadow-price vectors are exposed with the
\texttt{branch\_mu\_valid} flag; current native extraction does not certify
bounded branch-flow duals under congestion, so callers must check that flag
before interpreting branch duals.

\subsection{Reactive Power Optimization}
\label{sec:rpo}

\texttt{solve\_rpo} searches over discrete transformer tap and switchable-shunt
settings while calling AC OPF for the continuous sub-problem.  The implemented
workflow is outer-approximation / coordinate-search rather than a generic MINLP
branch-and-bound:
\begin{enumerate}
  \item sensitivity ranking by \(\pm1\) perturbations;
  \item ternary-search sweeps on ranked variables;
  \item coordinate-descent sweeps with OA lower-bound pruning;
  \item pairwise neighborhood polishing.
\end{enumerate}
The inner NLP uses the public AC OPF solver with parity IPM by default and
fallback enabled.
